\documentclass[10pt]{article}
\usepackage{template}
\NewsSetKeys{date=2026-10-14,vol=001,title={每日新闻·创刊号},author={ChatGPT},note={试刊版。仅收入重新核对过的消息；事实与评论分开处理。}}
\begin{document}
\begin{documents}
\newsitem{强化学习驱动的机器人发现不对称催化反应}{Nature Communications · 2026-10-13}{研究团队将深度强化学习与实验机器人结合，在数字孪生环境中训练策略，再迁移至真实的手性磷酸催化体系。30 天内系统做出超过 3000 次决策，发现四种高对映选择性转化，包括轴手性分子及三类联烯合成。\cite{Klochowicz2026}

\textbf{观察：}这项工作不仅是筛选既定反应条件，而是让算法在限定的反应空间内探索反应性、选择催化剂并扩展底物范围。但它不意味着自动化平台能在没有预先定义的化学空间时任意发现反应。}
\newsitem{NVIDIA Nova 推动 Rust DRM 基础设施}{Phoronix · 2026-10-13}{针对 Linux 7.4 的 DRM-Next 变更已送交审核，重点包括 Nova 驱动及其 Rust 基础设施。更新涉及 VRAM buddy 分配、TLB 刷新、MMU 页表类型、BAR1 映射的 CPU 访问和设备信息 ioctl。\cite{Larabel2026}

\textbf{观察：}这里的核心进展是系统层的资源管理、内存映射与驱动接口逐步成型。Rust 能减少部分内存安全风险，但不能自动消除硬件同步和驱动逻辑问题。合并目标不代表已进入稳定内核。}
\end{documents}
\bibliographystyle{plain}
\bibliography{refs}
\end{document}
